\scanpage{634}
\mspage{22}

du théorème d'adjonction de Quillen, d'où
l'adjonction de $u_!^W$ et de $u^*_{W}$, \unsure{cqfd|qed}.

\uline{Corollaire 2.4}. Soit \deleted{\unsure{$(M,W,C,F)$}} $W$ un localiseur
de $M$, tel qu'il existe un \deleted{\unsure{struct.}} triple
de Quillen spécial $(W,C,F)$ induisant $W$.
Alors le prédérivateur \deleted{défini} $\operatorname{Ho}_W$ \inserted{sur Cat} satisfait
aux axiomes Der 1 (décomposition), Der 2
(localisation), Der 3 (objet final), Der 4 bis
(\unsure{existence|existences} des $u_!^W$), Der 5 bis (calcul des
fibres des $u_!^W$), Der 6 bis (exactitude
à gauche, \deleted{\unsure{?}}\inserted{i.e. exactitude à g.} des carrés \inserted{$W$-}cocartésiens \unsure{dans|des} les $\operatorname{Ho}_W(1)$),
Der 6 (exactitude à droite \deleted{, des} carrés $W$-cartésiens).

On vient de prouver Der 4 bis, et Der. 5 bis
est \inserted{pratiquement} contenu dans le mode d'obtention de \deleted{\unsure{cet}}
$u_!$, compte tenu de la prop. 1.12, qui implique

\deleted{\uline{Corollaire \unsure{2.5}}. Soit $w\colon J\unsure{-|\longrightarrow} I$ une flèche de
Cat, supposons que l'on soit dans l'un des
trois cas suivants

a) $w$ est \unsure{coétale}

b) $J$ est discrète

c) $\forall\ i\in I$, $\deleted{\unsure{C}}J_{/i}$ est \unsure{bisomme|somme|réunion} de catégories ayant chacune un objet final.}
que pour \deleted{\unsure{$i'\text{ de }I'$}}
\deleted{\illegible{1}} $X$ dans $M(I)_c$, et $i'$ dans $I'$, l'image inverse
de $X$ dans $M({}_{\unsure{i'|i}\backslash}I)$ est dans $M({}_{i'\backslash}I)_c$ \ldots.
% archive sheet number: 633
